\section{Projektkontext und Relevanz}
\label{sec:01-02_project-context}

Der konkrete Projektkontext dieser Arbeit liegt in einem Transformationsvorhaben innerhalb der Informations- und Wissensorganisation eines großen Industrieunternehmens, in dem ein interner Wiki-Auftritt als zentrales Medium für Dokumentation, Methodenwissen und bereichsübergreifende Kollaboration etabliert werden soll. Die zuständige Transformationseinheit verfolgt das Ziel, komplexe Veränderungsinitiativen, neue Arbeitsweisen und organisationsweite Standards so aufzubereiten, dass Mitarbeitende aus unterschiedlichen Rollen und Standorten schnell auf relevante Informationen zugreifen und diese in ihren Arbeitsalltag integrieren können. Der betrachtete Wiki-Space fungiert dabei als zentrale Anlaufstelle für \Gls{playbook}, Richtlinien, Prozessbeschreibungen und Best Practices und übernimmt eine Schlüsselrolle für Wissensvermittlung, Onboarding und die kommunikative Begleitung organisationaler Veränderungen.

Eine vorangegangene Praxisanalyse des bestehenden Wiki-Auftritts hat jedoch gezeigt, dass die vorhandene Informationsarchitektur und \Gls{ux} den Anforderungen der heterogenen Nutzergruppen nur eingeschränkt gerecht werden. Nutzerbefragungen und qualitative Rückmeldungen verweisen auf unübersichtliche Navigationsstrukturen, schwer nachvollziehbare Seitenhierarchien, visuelle Inkonsistenzen sowie unklare Verantwortlichkeiten für Inhalte, was zu verlängerten Suchzeiten, Frustration und der verstärkten Nutzung alternativer Informationskanäle führt. In der Folge wird die Plattform im Alltag nur selektiv genutzt, während informelle Ablagen, E-Mail-Kommunikation oder persönliche Nachfragen weiterhin eine dominierende Rolle im Wissensaustausch spielen und damit genau jene \enquote{\Gls{singlesourceoftruth}} unterlaufen, die strategisch angestrebt wird \cite{402:Rosenfeld2015}.

Der Projektkontext ist zugleich eingebettet in eine umfassendere \gls{digitaleTransformation}, in deren Verlauf neue Methoden, \Gls{rollenmodelle} und Kollaborationsformen eingeführt und über digitale Plattformen skaliert werden sollen \cite{302:GSMA2024}. Studien zu Intranets und \Gls{enterprisesystem}e betonen, dass der Erfolg solcher Transformationsprogramme maßgeblich davon abhängt, ob zentrale Wissenssysteme intuitiv nutzbar, klar strukturiert und eng an reale Nutzungsszenarien gekoppelt sind, da sie als zentrale Kommunikations- und Lernmedien fungieren \cite{300:FraunhoferIESE2018,402:Rosenfeld2015}. Eine verbesserte \Gls{ux} interner Anwendungen wirkt sich nachweislich positiv auf Prozessqualität, Fehlerquote, Schulungsaufwand und Akzeptanz neuer Arbeitsweisen aus und wird von Forschungseinrichtungen wie Fraunhofer~IESE als wesentlicher Hebel zur Realisierung des wirtschaftlichen Nutzens digitaler Lösungen hervorgehoben \cite{300:FraunhoferIESE2018,301:Infragistics2015}.

Vor diesem Hintergrund verbindet die vorliegende Arbeit einen konkreten, praxisnahen Projektkontext mit einer übergeordneten wissenschaftlichen Fragestellung. Obwohl die Untersuchung aus einem spezifischen Transformationsprojekt hervorgeht, wird sie organisatorisch unabhängig durchgeführt und theoretisch so gerahmt, dass die gewonnenen Erkenntnisse und Gestaltungsprinzipien auf andere organisationale Wissens- und Kollaborationsplattformen übertragbar sind \cite{402:Rosenfeld2015}. Damit leistet die Arbeit einen Beitrag dazu, den Einfluss nutzerzentrierter \Gls{ux}- und \Gls{ia}-Strategien im organisationalen Kontext besser zu verstehen und wissenschaftlich fundierte, praxisrelevante Ansatzpunkte für die Gestaltung und Weiterentwicklung interner Plattformen bereitzustellen \cite{100:WhartonUX2025,300:FraunhoferIESE2018}. Eine detaillierte Darstellung des organisationalen Umfelds, der bestehenden Plattform, des aktuellen Stands von \Gls{ia} und \Gls{ux} sowie der beteiligten Nutzergruppen und Herausforderungen erfolgt in \hyperref[ch:03_pre-analysis]{Kapitel~3 \enquote{Voranalyse}} und bildet die empirische Grundlage für die im weiteren Verlauf entwickelten Forschungsfragen und Gestaltungsansätze.
